%% ── Chapter 3: Related Work ─────────────────────────────────
\chapter{Related Work}
\label{ch:related}

This chapter surveys four research streams underpinning \caesar{}.
For each stream we identify the key limitation that \caesar{} addresses.

\section{GAN-Based Intrusion Detection and Attack Synthesis}

Ring et al.\ \cite{ring2019} train a WGAN on CICIDS2017 for IDS data
augmentation. Lin et al.\ \cite{lin2018} introduce IDSGAN, using IDS
output as an adversarial signal to produce evasive traffic. Both
generators are \emph{unconditional} — they do not adapt to the current
defense state. \tagan{} closes this gap by conditioning $G_\phi$ on the
real-time defense embedding $\mathbf{d}_t \in \mathbb{R}^8$.

Xu et al.\ \cite{xu2019} introduce CTGAN for mixed-type tabular data.
\caesar{} extends this by treating the defense state as a continuous
conditioning vector rather than a categorical label.

\section{Reinforcement Learning for Cyber-Defense}

Elderman et al.\ \cite{elderman2017} model network defense as a
two-player zero-sum Markov game. Hammar and Stadler \cite{hammar2021}
use multi-agent RL for intrusion prevention but fix the attacker policy.
\caesar{} allows both attacker and defender to co-evolve, which is the
key novelty.

The Dueling Double-DQN architecture \cite{wang2016,hasselt2016} is the
dominant paradigm for discrete countermeasure selection
\cite{nguyen2021}. \adpn{} employs this and augments the reward signal
with co-evolutionary fitness terms.

\section{Adversarial ML in Network Security}

FGSM \cite{goodfellow2014b} and PGD \cite{madry2018} operate in the
gradient direction and often produce physically invalid network flows.
Pierazzi et al.\ \cite{pierazzi2020} distinguish feature-space from
problem-space attacks. \mgae{} bridges this gap via diffusion-guided
manifold projection.

Xiao et al.\ \cite{xiao2022} demonstrate that partial reverse diffusion
outperforms FGSM/PGD in image adversarial examples at equal distortion.
\mgae{} adapts this to tabular traffic data.

\section{Self-Healing Security Systems}

Coutinho et al.\ \cite{coutinho2014} propose a self-healing cloud
architecture using autonomic loops. Cheng et al.\ \cite{cheng2017}
apply RL to automated incident response but without generative attack
simulation. \caesar{}'s \arl{} unifies these streams.

\section{Positioning of \caesar{}}

\begin{table}[htbp]
\centering
\caption{Comparison of \caesar{} with representative related work.}
\label{tab:related}
\small
\begin{tabular}{@{}lccccc@{}}
\toprule
Work & \makecell{Def.-cond.\\attacker} &
       \makecell{Co-evol.\\loop} &
       \makecell{Proactive\\defense} &
       \makecell{Manifold\\evasion} &
       \makecell{Self-\\healing} \\
\midrule
IDSGAN \cite{lin2018}                & -- & -- & -- & -- & -- \\
Elderman et al.\ \cite{elderman2017} & -- & partial & -- & -- & -- \\
Hammar \& Stadler \cite{hammar2021}  & -- & -- & -- & -- & -- \\
Xiao et al.\ \cite{xiao2022}         & -- & -- & -- & partial & -- \\
\textbf{\caesar{} (ours)}            & \checkmark & \checkmark & \checkmark & \checkmark & \checkmark \\
\bottomrule
\end{tabular}
\end{table}

Table~\ref{tab:related} confirms no prior work combines all five
properties. The closest is Elderman et al., which includes partial
co-evolution but lacks conditioned generation, manifold evasion, and
self-healing.
